\begin{table}[ht]
\centering
\caption{Standardized Effect Sizes}
\label{tab:sde}
\small
\begin{tabular}{lcccccc}
\hline\hline
Outcome & $\hat{\beta}$ & SE & SD($Y$) & SDE & SE(SDE) & Classification \\
\hline
Charge rate (\%) & 12.919 & 3.487 & 4.075 & 2.167 & 0.585 & Large positive \\
log(Charges) & 0.526 & 0.182 & 0.721 & 0.500 & 0.173 & Large positive \\
Charge rate: violence & 18.555 & 4.660 & 6.515 & 1.931 & 0.485 & Large positive \\
Charge rate: sexual offences & 4.903 & 4.952 & 4.844 & 0.678 & 0.684 & Large positive \\
Charge rate: theft offences & 12.969 & 3.506 & 3.592 & 2.448 & 0.662 & Large positive \\
Charge rate: drug offences & 16.276 & 10.101 & 8.085 & 1.365 & 0.847 & Large positive \\
\hline\hline
\end{tabular}
\begin{minipage}{0.95\textwidth}
\vspace{3pt}
\small\textit{Notes:} This paper estimates the effect of police officer reductions on criminal justice quality in England and Wales, 2014--2021. Treatment: within-force variation in log officer FTE (continuous). Method: two-way fixed effects with force and year FE, clustered SEs. SDE = $\hat{\beta} \times \text{SD}(X) / \text{SD}(Y)$ for continuous treatment. N = 334 force-years (main); N = 628 (extended, row 2). Classification refers to magnitude of the point estimate, not statistical significance. Source: Home Office open data.
\end{minipage}
\end{table}

